% TOP_PROBLEM: {"id": 3700040, "title": "Defect relation for points in $\\mathbb P^2$", "category_id": 8, "category": {"id": 8, "name": "analysis", "display_name": "Analysis", "description": "Limits, continuity, calculus, and function theory.", "slug": "analysis", "order_index": 8, "created_at": "2026-07-31T15:26:25.671Z"}, "status": "open", "problem_number": "AMR-036-0040", "set_id": 13, "statusQualification": "Defines point proximity explicitly and uses the equivalent chordal sine of Fubini--Study distance."}
% FIRST_POSTED: 2026-10-01
% ENTRY_KIND: statement_only
% UnsolvedMath ID 3700040; source catalogue code AMR-036-0040
% !TeX program = lualatex
% Definition and sources only; this file is not an LLM solution attempt.
\documentclass[11pt,a4paper]{article}
\usepackage[margin=25mm]{geometry}
\usepackage{fontspec}
\setmainfont{lmroman10-regular.otf}[BoldFont=lmroman10-bold.otf,
 ItalicFont=lmroman10-italic.otf,BoldItalicFont=lmroman10-bolditalic.otf]
\usepackage{amsmath,amssymb,mathtools}
\usepackage{unicode-math}
\setmathfont{latinmodern-math.otf}
\AtBeginDocument{\let\setminus\smallsetminus}
\usepackage{xurl,hyperref}
\hypersetup{colorlinks=true,urlcolor=blue}
\setcounter{secnumdepth}{0}
\setlength{\parindent}{0pt}
\setlength{\parskip}{5pt}
\setlength{\emergencystretch}{3em}
\begin{document}
\section*{Defect relation for points in $\mathbb P^2$}\label{3700040}
{\small UnsolvedMath ID 3700040 · AMR-036-0040 · Analysis · AMR Open Problem Lists}
\subsection{Definitions and mathematical statement}
Let $f:\mathbb C\to\mathbb P^2(\mathbb C)$ be
a holomorphic curve with an entire homogeneous lift
$F=(F_0,F_1,F_2)$ having no common zero.
Require linear nondegeneracy: the functions $F_j$
are linearly independent over $\mathbb C$, equivalently
the image lies in no projective line.
Using the Euclidean norm on $\mathbb C^3$, define
\[
 T(r,f)=\frac1{2\pi}\int_0^{2\pi}
               \log\|F(re^{it})\|\,dt-\log\|F(0)\|.
\]
This characteristic is independent of the chosen
zero-free scalar change of lift.

For a point $a=[A]\in\mathbb P^2$, set
$d(f(z),a)=\|F(z)\wedge A\|/(\|F(z)\|\|A\|)$,
the sine of the Fubini--Study distance. Define
\[
 m(r,a,f)=\frac1{2\pi}\int_0^{2\pi}
                    \log\frac1{d(f(re^{it}),a)}\,dt,
 \qquad
 \delta(a,f)=\liminf_{r\to\infty}\frac{m(r,a,f)}{T(r,f)}.
\]
The logarithmic singularities at isolated preimages
are interpreted in the integral sense. Replacing
the sine distance by the ordinary Fubini--Study
distance changes the proximity by a bounded term,
and leaves this deficiency unchanged.

For every finite set of distinct points in
general position, meaning that no three are
collinear, must
$\sum_a\delta(a,f)\leq1$?
The same formulation for an infinite family uses
the supremum of its finite subsums. The targets
are points, of codimension two, rather than
hyperplanes. The question is the sharp simultaneous
defect bound for every linearly nondegenerate curve.
\subsection{Short English statement}
Is the total deficiency of points in general position at most one for every linearly nondegenerate entire curve in the projective plane?
\subsection{Sources}
\begin{itemize}
\item[S1] ULAM.AI and the UnsolvedMath Contributors, supplied problems.json, record 3700040 (AMR-036-0040), AMR Open Problem Lists. Catalogue identity and source transcription; CC BY 4.0.
\url{https://huggingface.co/datasets/ulamai/UnsolvedMath}.
\item[S2] Eremenko - My favorite unsolved problems, Defect relation conjecture. Original source indexed by AMR-036-0040.
\url{https://www.math.purdue.edu/~eremenko/uns1.html}.
\item[S3] A. Eremenko, Defect relation for targets of large codimension. Author-maintained primary problem note.
\url{https://www.math.purdue.edu/~eremenko/dvi/points.pdf}.
\end{itemize}
\end{document}
